\begin{textsample}{POS Dim 1 – vem\_ed\_32 – Score 19.00 – vem\_ed\_32/t171.txt}  \label{ex:f1_pos_034}
Osvaldo Succi em webinário da UNICollaboration Osvaldo Succi, coordenador dos Projetos Colaborativos Internacionais (PCIs) na equipe de Apoio à Internacionalização do Ensino Superior da Divisão de Extensão e Pesquisa no Ensino Superior (Depes) da Coordenadoria Geral de Ensino Superior de Graduação (CGESG) do Centro Paula Souza, ministrou webinário a convite da UNICollaboration em 31 de outubro de 2025.
A organização europeia tem o objetivo de \textbf{promover} a \textbf{integração} entre pesquisa e prática de Intercâmbios Virtuais.
O título da apresentação \textbf{foi} “The use of Artificial Intelligence by multilingual speakers and VE project coordinators” (O uso da Inteligência Artificial por falantes multilíngues e coordenadores de projetos de Intercâmbios Virtuais).
Entre os principais \textbf{pontos} desfavoráveis ao uso da IA por falantes multilíngues, Osvaldo Succi apontou limitações de \textbf{interação} \textbf{humana} e de \textbf{contexto} cultural; dependência excessiva e autonomia \textbf{reduzida}; \textbf{questões} éticas e de privacidade; barreiras técnicas e inequidade; além de qualidade e consistência.
Por outro lado, a IA ajuda os coordenadores de Intercâmbios Virtuais a aumentar a criatividade ao parear professores de \textbf{disciplinas} diferentes; pular perguntas básicas e fazer \textbf{questões} mais inspiradoras sobre as áreas \textbf{específicas} de \textbf{conhecimento} dos docentes; considerar diferentes tipos de Intercâmbio Virtual (em especial os inter e transdisciplinares); manter melhores registros das reuniões; lidar com \textbf{tarefas} repetitivas, porém importantes, como detalhar instruções para os alunos e, assim, liberar tempo para os professores parceiros conversarem e \textbf{fortalecerem} seus laços de relacionamento.
Enquanto a inteligência \textbf{humana} \textbf{é} melhor para eticamente avaliar soluções, aplicando sabedoria e \textbf{contexto}, a \textbf{artificial} auxilia a estimular cenários e sugerir soluções orientadas por dados.
Se a IA proporciona acesso imediato a quantidades enormes de informação, por outro lado a qualidade \textbf{humana} da curiosidade, de desaprender e reaprender \textbf{é} única.
O importante \textbf{é} \textbf{que} sempre haja a mediação \textbf{humana} para somar e ponderar o melhor de \textbf{cada} uma das inteligências, com ética e espírito \textbf{crítico}.

% matched lemmas: artificial, cada, conhecimento, contexto, crítico, disciplina, específico, fortalecer, humano, integração, interação, ponto, promover, que, questão, reduzir, ser, tarefa
\end{textsample}
